\clearpage
\hypertarget{gaussix-sources}{}
\section*{Sources and textual notes}

\subsection*{Historical edition and contributors}

\noindent Carl Friedrich Gauss, \textit{Werke. Neunter Band}. Issued by the Königliche Gesellschaft der Wissenschaften zu Göttingen. In commission with B. G. Teubner, Leipzig, 1903. Felix Klein served as an editor of the collected works and signed the closing \textit{Bemerkungen zum neunten Bande} as F. Klein. \href{https://gdz.sub.uni-goettingen.de/id/PPN23601515X}{Göttinger Digitalisierungszentrum catalogue}. \href{https://lbezone.hkust.edu.hk/rse/?p=13258}{HKUST Library catalogue}.

Carl Friedrich Gauss's mathematical papers, manuscripts, and correspondence form the body of the volume. In his remarks on p.~521, Klein credits L. Krüger with preparing Volume IX and A. Börsch with helping to examine the material. M. Brendel directed the general editorial work, as in Volume VIII. The historical remarks retain their signatures, including those of Krüger and the joint signatures of Krüger and Börsch. Hartmann is credited in the station observations. H. C. Schumacher wrote the introduction to the first published Hanover report. The remarks on station adjustments quote Goldschmidt's \textit{Geodätischer Calcül nach Gauss}. Lieutenant J. Gauss devised the Hanover triangulation map. Krüger's accompanying historical note identifies the editorial corrections to that map.

The historical editors identify additions in square brackets or in the remarks, and enclose excerpts of notes and letters by other authors in braces. Klein's account of these conventions is on p.~521.

The title lettering and arrangement follow the 1903 Göttingen copy. The unsigned vignette is documented on the \href{https://gdz.sub.uni-goettingen.de/id/PPN235993352}{1863 title leaf of \textit{Werke} I}. The image here is the matching vignette from the second printing of \textit{Werke} I (1870), \href{https://commons.wikimedia.org/wiki/File:Gauss,_Carl_Friedrich_-_Vignette_-_von_Werke_(1870)_(page_9_crop).jpg}{Wikimedia Commons, revision 1179960952}. \href{https://creativecommons.org/publicdomain/mark/1.0/}{Public Domain Mark 1.0}. The historical text, formulas, tables, and mathematical diagrams are set in editable LaTeX with modern fonts and spacing. The drawings follow the printed labels and connections, with approximate geometry where noted below. The edition reproduces the historical preliminaries and their three blank versos. The modern digitizer's cover, library stamps, and manuscript accession marks belong to individual copies.

\subsection*{Contents of Volume IX}

\noindent Page numbers below are those of the 1903 edition, including the contextual positions of unnumbered leaves. The final blank leaf and publisher's advertisement follow p.~528. The historical table of contents retains its original page references.

\paragraph{The latitude difference between Göttingen and Altona (pp.~1--64).}
\textit{Bestimmung des Breitenunterschiedes zwischen den Sternwarten von Göttingen und Altona durch Beobachtungen am Ramsdenschen Zenithsector} contains the title, contents and introduction, the catalogue of 43 stars, and the observation series. The results give the weighting formulas, collimation errors, residuals, statistical reductions, direct and reflected observations, and star comparisons. The Seeberg observations comprise sixty results. The \textit{Zusatz} lists twenty-five survey localities, and the \textit{Anzeige} includes historical remarks.

\paragraph{The Earth Ellipsoid and the Geodesic Line (pp.~65--98).}
The title leaf and its blank verso precede three posthumous articles on the ellipsoid with their numerical calculations and remarks, followed by Articles [1]--[13] on the geodesic line. These contain auxiliary tables, two 61-row tables, worked calculations, mathematical diagrams, and the historical remarks on pp.~96--98.

\paragraph{Latitude and spherical trigonometry (pp.~99--104).}
The letter to J. J. Baeyer on the variation of latitude with height occupies pp.~99--102, with a diagram and historical remarks. Two further posthumous notes concern the angles of spherical triangles and points on a great circle, with a diagram and remarks signed by Krüger and Börsch.

\paragraph{Conformal Double Projection (pp.~105--122).}
\textit{Conforme Doppelprojection des Sphäroids auf die Kugel und die Ebene} opens with a title leaf and blank verso. Articles [1]--[10] on transferring the ellipsoid to the sphere include a numerical table and historical remarks. The stereographic projection of the sphere onto the plane follows, with its remarks.

\paragraph{Mercator Projection and the stereographic representation of the spheroid (pp.~123--134).}
Articles [1]--[8] on Mercator projection include a numerical table and remarks signed by Krüger and Börsch. The stereographic representation of the spheroid in the plane contains formulas and the historical remarks on pp.~133--134.

\paragraph{Conformal Projection of the Spheroid onto a Cone (pp.~135--140).}
The title leaf and blank verso precede the posthumous text, its formulas and two numerical tables, and the historical remarks.

\paragraph{Conformal Mapping of the Spheroid in the Plane (pp.~141--204).}
The title leaf and blank verso precede the posthumous Articles [1]--[32] and historical remarks. These include coordinate transformations, logarithmic and curvature formulas, numerical tables, and the worked calculation for Varel.

\paragraph{On the Formulas for the Hanover Land Survey (pp.~205--218).}
The letters to Schumacher are dated 18 April 1830, 30 April 1830, 17 May 1831, 25 June 1831, and 9 December 1838. They include auxiliary tables and calculations for Neuwerk. The historical remarks are signed by Krüger and Börsch.

\paragraph{Trigonometrical Point Determination (pp.~219--242).}
The title leaf and blank verso precede Articles [1]--[9] and historical remarks with diagrams, tables, the Goldschmidt quotation, and the credit to Krüger.

\paragraph{Adjustment of Simple Figures (pp.~243--260).}
The title leaf and blank verso precede Articles [1]--[7], the letters of 11 February 1824 and 19 January 1840, and remarks signed by Krüger. The articles include diagrams, formulas, and tables.

\paragraph{Station Adjustments (pp.~261--293).}
The title leaf and blank verso precede the observations and normal equations for Zeven, Brillit, Wilsede, and Windberg, with the historical credit to Hartmann. They are followed by the Breithorn observations of 13--16 September 1822, the letter to Gerling of 26 December 1823, and those to Schumacher of 22 December 1827 and 7 January 1828. The tables, triangle drawing, and historical notes are followed by Krüger's remarks, which quote Goldschmidt's \textit{Geodätischer Calcül nach Gauss}.

\subsubsection*{On Network Adjustment (pp.~294--328)}
The title leaf and adjoining blanks precede the condition-count formulas [1], the 33 connection-list entries and network with 75 edges [2], and the 51 observed-triangle groups with their concluding relations [3]. Article [4] contains 51 numbered equations. Article [5] contains the introduction and twelve worked cases I--XII. Article [6] gives the relation table, correction expressions, and twelve normal equations. Article [7] gives the 150 direction corrections and concluding text. Article [8] gives 51 triangle groups and the Hamburg--Hohenhorn connection table. Article [9] gives 150 azimuth and reduction rows. Article [10] comprises excerpts of Gauss's letters to Olbers of 2 November 1823, 6 July 1824, and 14 May 1826, and to Gerling of 2 May 1837, 5 June 1838, and 14 November 1838. It includes the historical footnotes, the six-point triangulation diagram and quadrilateral ring diagram, and Krüger's historical remarks.

\paragraph{Ring of Triangles Around Oldenburg (pp.~329--342).}
The title leaf, blank verso, and opening paragraph precede 84 preliminary table rows in 21 station groups, a triangulation diagram with 21 numbered vertices and 42 edges, the 22 angle-sum errors, condition examples, 22 normal equations and correlates, and the discussion of corrections. The triangle overview gives 63 rows of angles and side logarithms. The coordinate table has 23 coordinate rows and two difference rows. The second adjustment gives the formulas with six direction relations, 77 coefficient pairs, and their explanation. Krüger's historical remarks quote Gauss's letter to v.~Schrenck of 3 August 1835. The historical adjustment itself remains unfinished.

\paragraph{On the Hanover Triangulation (pp.~343--434 and the foldout).}
The title leaf and blank verso precede the correspondence. Section [1] comprises eight excerpts of letters to Schumacher between 5 July 1816 and 14 January 1829 and three historical footnotes. Section [2] comprises eight excerpts of letters to Bessel of 26 December 1821, 15 November 1822, 5 November 1823, 20 November 1824, 12 March 1826, 20 November 1826, 1 April 1827, and 9 April 1830, with three historical footnotes. Section [3] contains the letter to Bohnenberger of 16 November 1823 and its three historical footnotes. Section [4] contains twelve excerpts of letters to Olbers dated 13 January 1821, 18 April 1822, 4 July 1824, 8 July 1824, 19 February 1825, 25 February 1825, 9 October 1825, 2 April 1826, 14 January 1827, 1 March 1827, 14 June 1830, and 2 September 1837. Three historical footnotes, the Bremen--W.--Brüttendorf--Bottel triangulation diagram, and the 57-row \textit{Relative Höhen} table accompany these excerpts. Section [5] contains eleven excerpts of letters to Gerling dated 5 October 1821, 21 February 1822, 7 November 1822, 27 July 1823, 11 August 1823, 1 September 1823, 5 September 1823, 3 October 1823, 19 July 1827, 12 September 1838, and 14 November 1838, with four heliotrope signalling diagrams, three optical diagrams, a coordinate table, a historical footnote, and remarks signed by Krüger.

Section [II.] contains the published reports of 1821 and 1822, Schumacher's introduction, the coordinate table, and the historical remarks. Section [III.1] gives the plan of the work, its introduction, and the beginning of its first section. Section [III.2] contains excerpts from the reports of 7 January 1822, 31 January 1823, 16 February 1825, 21 November 1827, and 26 June 1828, the historical report of 8 February 1838, and excerpts dated 5 July 1840 and December 1844. Section [III.3] contains the complete thirty-three-row coordinate table and Krüger's historical remarks. The editable foldout depicts the triangulation networks, geographical outlines, legends, and scale.

\paragraph{Height Measurements (pp.~435--458).}
The title leaf and blank verso precede three texts on refraction [1]--[3], including the letters to Schumacher of 28 November 1824 and 27 December 1846 and the historical editorial note. The posthumous text on terrestrial refraction and height adjustment includes 71 error equations, 31 normal equations, station corrections, and Krüger's remarks. Three tables accompany the instructions for barometric height measurement.

\paragraph{The Heliotrope (pp.~459--484).}
The title leaf and blank verso precede the article in the \textit{Göttingische gelehrte Anzeigen} of 9 August 1821, with an optical diagram, and the contribution to the \textit{Astronomisches Jahrbuch} submitted on 26 December 1821. The letter to Olbers is dated 1 July 1821. The four letters to Schumacher are dated 11 July 1821, 8 November 1821, 10 July 1822, and 15 January 1827, with their notes and drawing. The article on adjusting the heliotrope includes a thirteen-figure plate. Two posthumous notes include formulas and a geometric drawing. Krüger's historical remarks give a diagram of the first construction.

\paragraph{Measurement Errors (pp.~485--506).}
The title leaf and blank verso precede the text on surveying instruments, the letters on measurement errors to Olbers (July 1825), Schumacher (14 August 1825), Bessel (29 October 1843 and 15 August 1844), and Gerling (17 April 1844), and the observation calculations. Krüger signed the historical remarks.

\paragraph{Reduction of Oblique Angles to the Horizon. Observations by Repetition. (pp.~507--520).}
The title leaf and blank verso precede three notes on the reduction of oblique angles to the horizon, a worked example, the calculation scheme, and the historical remarks. The texts on observations by repetition include the letter to Bessel of 27 January 1819, the recurrence and error formulas, and the historical remarks signed by Krüger.

\paragraph{Remarks on the Ninth Volume and Contents (pp.~521--528).}
F. Klein signed the historical remarks on the ninth volume. The original contents, final blank leaf, and publisher's advertisement close the volume.

\subsection*{Source copies and map}

The translation is based on the original-language AI transcription of Volume IX dated 11 November 2026, \texttt{gauss-werke-ix-1-528-latex.zip}.

The digitized Göttingen copy of the 1903 edition (PPN23601515X) is the principal source. The Google Books copy at the Internet Archive and the HKUST Library copy are comparison witnesses. SUB Göttingen separately digitized the unfolded Hanover map, \textit{Übersicht der gemessenen Dreieckssysteme}. The Göttingen volume PDF omits the map. The Google volume shows folded strips.

\noindent\href{https://gdz.sub.uni-goettingen.de/download/pdf/PPN23601515X/PPN23601515X.pdf}{Principal source scan, Göttinger Digitalisierungszentrum}. \href{https://archive.org/details/carlfriedrichga01gausgoog9}{Google Books copy, Internet Archive}. \href{https://lbezone.hkust.edu.hk/rse/?p=13258}{Comparison copy, HKUST Library}. \href{https://webdoc.sub.gwdg.de/ebook/e/2005/gausscd/html/kapitel_landv_karte.htm}{Hanover foldout map, SUB Göttingen}.

The blank leaves within pp.~1--528 occupy pp.~2, 4, 66, 106, 136, 142, 220, 244, 262, 294, 296, 330, 344, 436, 460, 486, 506, and 508. The historical title leaves, final unnumbered blank leaf, and publisher's advertisement are included.

\subsection*{Textual and graphical readings}

The quoted spellings are from the historical copies. The notes also describe the rendition of mathematical signs and figures in this edition.

The HKUST copy clarifies faint delimiters or small notation on pp.~53--54, 57--58, 60, 64, 69--72, 76--77, 80, 82, 86, 98, 102, 109--112, 119, 125, 129, 132, 138, 144, 152--162, 171, 195--199, 201--202, and 212. It also supports the network and connection list on pp.~298--299, both \(\tfrac13\,\text{Excess}\) table headings on pp.~304--305, and the damaged letter in \textit{20-maliger} on p.~441.

\noindent\href{https://lbezone.hkust.edu.hk/obj/4/o/b761099/v9/pg00220.jpg}{HKUST, p.~212}. \href{https://lbezone.hkust.edu.hk/obj/4/o/b761099/v9/pg00312.jpg}{HKUST, p.~304}. \href{https://lbezone.hkust.edu.hk/obj/4/o/b761099/v9/pg00313.jpg}{HKUST, p.~305}. \href{https://lbezone.hkust.edu.hk/obj/4/o/b761099/v9/pg00451.jpg}{HKUST, p.~441}.

The two low marks in the running head on page 6 and the trailing mark
after 54,24 in the page 18 observation table follow the principal copy. Their functions are uncertain. The star designation on page 9 uses the
same qualified kappa variant discussed below. The colon at \textit{Altona} on page 12 and \textit{Jun} without a period on page 15 follow the principal text.

The principal copy's separate point after \(\beta'\) in the weighting
formula on page 30 is retained. Its function is uncertain. The differing
latitude values on pages 38 and 39 follow the printed text without correction.
The signature on page 35 is supplied from the Google and HKUST copies.

The separate point before \textit{den} on page 51 follows the principal copy, with uncertain function. The gathering signature on that page follows the Google copy. The compact cross-shaped Greek designation in star row 8 is represented by a kappa variant. The exact historical typeform is uncertain. Numerical and geographical spellings in the Seeberg and survey tables follow the printed text, including unreduced seconds.

Faint citation delimiters on pages 64 and 71, small digits on page 69,
formula delimiters and the outer square on page 70, and a faint
parenthesis on page 72 follow the clearer HKUST copy. The separate point
before the final logarithm on page 69 follows the principal copy. Its function remains uncertain. The auxiliary table on page 77 preserves
the printed blank cells and abbreviated values. Diagrams on pages 73
and 75 preserve their labelled connections with approximate geometry.

The clipped rightmost column of the table on page 82 follows the HKUST
copy. The flattening values on pages 82 and 84 differ from those on
pages 83 and 87, as printed. The separate points on pages 84, 88, 93 and
94 and the raised point in the formula ending on page 91 are retained
without assigning them an additional mathematical function. A faint
closing parenthesis on page 86 and the fourth power on page 98 follow
the clearer HKUST copy. The diagrams on pages 100 and 103 use approximate
geometry while preserving their labelled connections and crossings.

The page range 30--32 cited on page 102 and small notation on pages
109--112 follow the clearer HKUST copy. The separate point before the
denominator on page 113 follows the principal copy. Its function is uncertain. The vertical stroke after cos on page 119 is present in the
principal, Google and HKUST copies. Its function remains uncertain.
Unclassified peripheral ink on pages 112, 117--118 and 120--121 is not
represented as additional mathematical characters.

The lowered point beneath 2345.09 on page 126 is retained separately
from the digits. Its function is uncertain. The short stroke near 8504.24
is unclassified. A faint closing parenthesis on page 129 follows the
clearer HKUST copy. That copy also clarifies the denominator 45360 and
delimiters on page 132 and the series on page 138.
On page 139 the gaps in 6,709 8436 473 and 1180 553 follow the principal
copy. Google has points at those positions. The separate point before
860 in the lower table follows the principal and is absent from Google.
The separate points on pages 138--140 are retained with uncertain functions.
Unclassified peripheral ink on pages 132, 137 and 140 is not represented
as additional mathematical characters.

Faint closing parentheses on page 144 follow the clearer HKUST copy. Those in the first \(\log\) hyp \(D\) formula on page 147 follow Google.
On page 160 the point after 89,437 follows the principal copy and is absent
from Google and HKUST. The two marks before 109 on page 161 and the
raised quote-like mark before 89 on page 162 are retained with uncertain
functions. The table value 788 on page 162 and denominator 5040 on
page 171 follow the clearer comparison witnesses. The separate points
on pages 163, 167, 169--171 follow the principal copy. Their functions are uncertain. Unclassified peripheral ink on pages 144, 147--150 and
166 is not represented as additional mathematical characters.

On page 195 the comma after \textit{zugefügt worden} follows the principal copy. Google and HKUST have a semicolon. On page 197 the period after the first
displayed 3,418 9411.803 follows the principal copy. The comparison copies have a comma.
The small coefficient on page 195 and cramped notation on pages 196--199
and 201--202 are supported by the clearer HKUST copy. The unclassified
isolated marks on pages 192 and 198 and in the otherwise blank degree cell
of the page 210 table are not represented as additional characters.
The separate points in the calculation field on page 203, after Es and
beside the first longitude equation on page 208, and above the minus
in formula 6) on page 209 follow the principal copy. Their functions are uncertain.

On page 212 the exponent 2 in the denominator of \(F\) follows the clearer
HKUST copy. The upper and lower marks between 434 and 2945 in \(k\)
follow the principal copy. Their functions remain uncertain. The separate
short stroke after the comma following \textit{gebrauchen} on page 225 follows
the principal copy and is absent from Google. Its textual status is uncertain.

On page 233 the principal semicolon after \textit{scheint} differs from the counterpart comma. The raised comma-like mark by the title on page 239 has uncertain function.
On page 242 the principal period after \textit{sein} differs from the counterpart colon.
Its point before \textit{Rechnungen} and the raised point before \textit{eine} on page 249 are
retained with uncertain function. The obscured text on page 250 of the Google copy follows the principal copy.
The principal scan clips the right-hand part of the system on page 270. The complete Google copy supplies that portion and the sideways orientation.
Unreadable areas of Google page 254 follow the principal copy. Separate marks
on pages 254--257, 260, 263--264 and 268--269 follow the principal copy. Their functions remain uncertain. The repeated e headings on page 265 are retained.

The separate marks beside the formulas on pages 272--273, the folios
on pages 274 and 276, the value 253 on page 276, the footnote on
page 278 and the two prose/display positions on page 281 follow the
principal copy. Their functions and physical origins are uncertain.
The closing rule on page 277 occurs in the Google copy alone.
On page 286 the first item has a comma after I in the principal copy. The comparison copy has a point. The inferior logarithmic marks on
page 287 are represented by \(n\). Their exact printed type is uncertain.
The separate marks on pages 292 and 297 and in the Wilsede row on
page 298 follow the principal copy. Their functions are uncertain.
The first Wilsede leader is comma-shaped in the principal copy and
point-shaped in HKUST. On page 300 the triangle 20 sum has a point
in 0.545, whereas Google has a comma. The separate point in triangle 8
and the dash below the middle table follow the principal positions.
The upper point above the comma in triangle 48 on page 301 follows
the principal copy. Isolated marks on pages 283, 293 and 301 occur in the comparison copies alone. Unidentified peripheral ink on pages 282--285
is not represented as additional characters. The page 299 network
preserves the printed incidences and labels with approximate geometry.

On p.~444 the principal copy reads \textit{Nach diesem}, referring to the handbook, and puts the subsequent reference to that handbook in parentheses. The rotation abbreviations on pp.~481--482 retain the German initials L (\textit{links}), O (\textit{oben}), R (\textit{rechts}), and U (\textit{unten}).

The signature on page 465 and the terminal rule on page 466 follow the
Google comparison copy and are absent from the principal scan.
The isolated points on pages 473, 474 and 480 follow the principal copy. Their function is uncertain. The plate labels on page 477 retain the
printed J and lower-case c, differing from the prose references I and CC.
The principal copy has a semicolon after Müller on page 489, where Google has a comma. The raised point on page 495 follows the principal copy and is absent from Google.
The terminal point after \textit{Bremervörde} in the page 492 diagram is uncertain
in the principal copy and dot-like in Google. The irregular ink attached
to the readable word \textit{Ich} on page 495 is not reproduced as an additional sign.
The separate points before \textit{Fernrohr} on page 501 and \textit{Fehler} on page 515
follow the principal copy and are absent from Google. Their function is uncertain.
The signature IX. / 64 on page 505 follows Google and is absent from the
principal scan. Faint parentheses on page 512 follow the readable comparison
copy. The bracket-like mark between the equals sign and cube root in the
definition of \(\mu\) on page 513 is represented by \texttt{]}. Its function is uncertain.
Irregular peripheral ink near \textit{Aus} and the final exponent on page 516 is
not represented as additional text.
On page 519 the exponent 2 in the elimination equation and exponent 4
in the reciprocal-weight numerator follow the principal copy. The Google superscripts are filled. The date 18 September 1822 on page 526 follows
the principal copy. Google appears to read 19. The separate point after
\textit{Auszug} on page 526 is retained. Unidentified peripheral ink near the
heading on page 521, \textit{Azimuth} on page 524 and \textit{an Olbers} on page 526 is
not represented as additional characters. On page 528 the principal copy prints \textit{Uber Messungsfehler} with unaccented \textit{Uber}, and gives \textit{Messungsfehler} and \textit{Repetitionsbeobachtungen} without colons. The extra point in
the 1822 Juli 10 leader and comma-shaped mark before heading B. in the
advertisement are retained with uncertain functions. The advertisement follows the principal final leaf. It has no counterpart in the Google copy.

The unfolded Hanover map follows the separate SUB Göttingen copy. Its geometry, stroke weights, and colour-band offsets approximate the printed drawing. Overprinted contour centre lines at Pugelatz, Fuhrenberg, and the ray to Inselsberg cannot be distinguished throughout. The obscured portions remain open.

The quotation from Gauss's letter to v.~Schrenck of 3 August 1835 also appears in the \textit{General-Bericht über die mitteleuropäische Gradmessung für das Jahr 1865}. Berlin, Georg Reimer, 1866, pp.~23--24.\par
\noindent\href{https://books.google.com/books?id=VBNDAAAAcAAJ\&pg=PA24}{Earlier printing of the letter, Google Books}.

